% ============================================================================
% SOLUCIONES - MÓDULO 9 (PROYECTO FINAL: RETOS DE EXTENSIÓN)
% ============================================================================

\begin{solucion}{p-1}
La recompra se mide sobre los clientes \emph{activos} (los que compraron al
menos una vez). Para los días entre compras ordenamos las compras de cada
cliente por fecha y restamos la fecha anterior con \codigo{lag()}.

\begin{rcode}
library(tidyverse)
transacciones <- read_csv("datasets/transacciones.csv",
                          show_col_types = FALSE)
\end{rcode}

Al cargar \paquete{tidyverse} aparece su mensaje de bienvenida (no lo
copiamos aquí). Ahora los tres cálculos:

\begin{rcode}
# (a) Porcentaje de clientes activos con 2 o más compras
compras_cliente <- transacciones %>% count(cliente_id, name = "compras")
compras_cliente %>%
  summarise(activos = n(),
            recompran = sum(compras >= 2),
            tasa_recompra = round(100 * mean(compras >= 2), 1))

# (b) Días entre compras consecutivas del mismo cliente
intervalos <- transacciones %>%
  arrange(cliente_id, fecha) %>%
  group_by(cliente_id) %>%
  mutate(dias_desde_anterior = as.numeric(fecha - lag(fecha))) %>%
  ungroup() %>%
  filter(!is.na(dias_desde_anterior))

quantile(intervalos$dias_desde_anterior, c(0.25, 0.5, 0.75))

# (c) Clientes con al menos una recompra en 30 días o menos
intervalos %>%
  group_by(cliente_id) %>%
  summarise(recompra_30 = any(dias_desde_anterior <= 30)) %>%
  summarise(clientes = sum(recompra_30))
\end{rcode}

\begin{rsalida}
# A tibble: 1 x 3
  activos recompran tasa_recompra
    <int>     <int>         <dbl>
1     198       192            97
25% 50% 75% 
 19  41  82 
# A tibble: 1 x 1
  clientes
     <int>
1      140
\end{rsalida}

\textbf{Interpretación:} 192 de los 198 clientes activos (97\%) compraron
dos veces o más, así que ``conseguir la segunda compra'' no es el problema de
TechRetail. La mediana entre compras es de 41 días (la mitad de las
recompras ocurre entre 19 y 82 días) y 140 clientes volvieron alguna vez en
30 días o menos. Una campaña de segunda compra aportaría poco; es más útil
una campaña que \emph{acorte} el intervalo (por ejemplo, un cupón válido 30
días) y medir si la mediana baja de 41 días.
\end{solucion}

\begin{solucion}{p-2}
Unimos las ventas por sucursal con \archivo{tiendas.csv}, dividimos entre
m\textsuperscript{2} y empleados y usamos \codigo{min\_rank(desc())} para
los rankings.

\begin{rcode}
library(tidyverse)
transacciones <- read_csv("datasets/transacciones.csv",
                          show_col_types = FALSE)
tiendas <- read_csv("datasets/tiendas.csv", show_col_types = FALSE)

productividad <- transacciones %>%
  group_by(tienda_id) %>%
  summarise(ventas = sum(total_transaccion), .groups = "drop") %>%
  left_join(tiendas, by = "tienda_id") %>%
  mutate(ventas_m2       = ventas / tamano_m2,
         ventas_empleado = ventas / empleados,
         rank_ventas     = min_rank(desc(ventas)),
         rank_m2         = min_rank(desc(ventas_m2)),
         rank_empleado   = min_rank(desc(ventas_empleado)),
         cambio          = rank_ventas - rank_m2) %>%
  arrange(rank_m2)

productividad %>%
  transmute(tienda = str_remove(nombre_tienda, "Sucursal "),
            m2 = tamano_m2, ventas = round(ventas),
            por_m2 = round(ventas_m2), por_empleado = round(ventas_empleado),
            rank_ventas, rank_m2, cambio)
\end{rcode}

\begin{rsalida}
# A tibble: 10 x 8
   tienda     m2 ventas por_m2 por_empleado rank_ventas rank_m2 cambio
   <chr>   <dbl>  <dbl>  <dbl>        <dbl>       <int>   <int>  <int>
 1 Express   300 102791    343         8566           4       1      3
 2 Outlet    350 114377    327         7625           1       2     -1
 3 Este      400  88420    221         4912           7       3      4
 4 Norte     450  89024    198         4451           6       4      2
 5 Centro    500  88414    177         3537           8       5      3
 6 Sur       600 105305    176         3510           2       6     -4
 7 Plaza B   480  80236    167         3821          10       7      3
 8 Oeste     550  84672    154         3849           9       8      1
 9 Plaza A   700 105216    150         3006           3       9     -6
10 Premium   800 101538    127         2538           5      10     -5
\end{rsalida}

\begin{rcode}
color_resalte <- "#eb6834"; color_gris <- "#b5b3ad"
productividad %>%
  mutate(mejor = rank_m2 == 1) %>%
  ggplot(aes(x = ventas_m2, y = fct_reorder(nombre_tienda, ventas_m2),
             fill = mejor)) +
  geom_col(width = 0.7, show.legend = FALSE) +
  scale_fill_manual(values = c(`TRUE` = color_resalte,
                               `FALSE` = color_gris)) +
  scale_x_continuous(labels = scales::dollar) +
  labs(title = "Ventas por metro cuadrado, 2023",
       x = "Ventas por m²", y = NULL) +
  theme_minimal() +
  theme(panel.grid.minor = element_blank())
\end{rcode}

\textbf{Interpretación:} la Express (300~m\textsuperscript{2}) es la más
productiva por metro cuadrado (\$343) y por empleado (\$8,566), seguida de la
Outlet. La Plaza A es la que más cae: tercera en ventas totales y novena por
m\textsuperscript{2} (6 lugares), y la Premium queda última por
m\textsuperscript{2} (\$127). Si la dirección piensa en expansión, los
formatos chicos producen casi el triple por metro cuadrado que los grandes;
antes de decidir habría que agregar el costo de renta por
m\textsuperscript{2}, que no está en los datos.
\end{solucion}

\begin{solucion}{p-3}
La venta diaria promedio es el total de unidades de 2023 entre 365 días. La
cobertura dice cuántos días alcanza el stock actual a ese ritmo.

\begin{rcode}
library(tidyverse)
transacciones <- read_csv("datasets/transacciones.csv",
                          show_col_types = FALSE)
productos <- read_csv("datasets/productos.csv", show_col_types = FALSE)

alerta <- transacciones %>%
  group_by(producto_id) %>%
  summarise(unidades_2023 = sum(cantidad), .groups = "drop") %>%
  right_join(productos, by = "producto_id") %>%
  filter(activo) %>%
  mutate(unidades_2023   = replace_na(unidades_2023, 0),
         venta_diaria    = unidades_2023 / 365,
         dias_cobertura  = stock_actual / venta_diaria,
         bajo_minimo     = stock_actual < stock_minimo,
         alerta          = dias_cobertura < 30 | bajo_minimo) %>%
  filter(alerta) %>%
  arrange(dias_cobertura) %>%
  transmute(producto_id, stock = stock_actual, minimo = stock_minimo,
            diaria = round(venta_diaria, 2),
            cobertura = round(dias_cobertura), bajo_minimo)

alerta

# Contexto: distribución de la cobertura de todos los activos
productos %>%
  filter(activo) %>%
  left_join(count(transacciones, producto_id, wt = cantidad,
                  name = "unidades"), by = "producto_id") %>%
  mutate(dias_cobertura = stock_actual / (unidades / 365)) %>%
  pull(dias_cobertura) %>%
  quantile(c(0, 0.25, 0.5)) %>%
  round()
\end{rcode}

\begin{rsalida}
# A tibble: 3 x 6
  producto_id stock minimo diaria cobertura bajo_minimo
        <dbl> <dbl>  <dbl>  <dbl>     <dbl> <lgl>      
1           3    22     26   0.18       125 TRUE       
2           6    36     38   0.22       160 TRUE       
3          33    42     45   0.21       199 TRUE       
  0%  25%  50% 
 125  927 1717 
\end{rsalida}

\textbf{Interpretación:} ningún producto activo tiene menos de 30 días de
cobertura: el que menos tiene alcanza para 125 días. Las 3 alertas se deben
a que el stock está por debajo del mínimo definido (productos 3, 6 y 33),
aunque al ritmo de venta actual (unas 0.2 unidades diarias) les alcanza para más
de cuatro meses. El hallazgo importante es el contrario: la mediana de
cobertura es de 1,717 días (más de cuatro años de inventario). Con 1,000
transacciones al año, TechRetail parece tener \emph{sobreinventario}, y los
\codigo{stock\_minimo} están desalineados con la demanda real. Recomendación:
recalcular mínimos a partir de la venta diaria (por ejemplo, 30 días de
cobertura) y revisar si el stock se registra en unidades o en cajas.
\end{solucion}

\begin{solucion}{p-4}
La cohorte es el mes de la primera compra. Para cada compra calculamos
cuántos meses pasaron desde ese mes (\codigo{meses\_desde}) y contamos
clientes distintos por cohorte y mes.

\begin{rcode}
library(tidyverse)
transacciones <- read_csv("datasets/transacciones.csv",
                          show_col_types = FALSE)

compras <- transacciones %>%
  mutate(mes_compra = floor_date(fecha, "month")) %>%
  group_by(cliente_id) %>%
  mutate(cohorte = min(mes_compra)) %>%
  ungroup() %>%
  mutate(meses_desde = (year(mes_compra) - year(cohorte)) * 12 +
                        month(mes_compra) - month(cohorte))

# Tamaño de cada cohorte
compras %>% distinct(cliente_id, cohorte) %>% count(cohorte) %>%
  mutate(cohorte = format(cohorte, "%Y-%m"))

matriz <- compras %>%
  distinct(cliente_id, cohorte, meses_desde) %>%
  count(cohorte, meses_desde, name = "clientes") %>%
  group_by(cohorte) %>%
  mutate(pct = round(100 * clientes / clientes[meses_desde == 0])) %>%
  ungroup()

# Primer trimestre: % de la cohorte que compró 1, 2 y 3 meses después
matriz %>%
  filter(cohorte <= as.Date("2023-03-01"), meses_desde <= 3) %>%
  mutate(cohorte = format(cohorte, "%Y-%m")) %>%
  select(-clientes) %>%
  pivot_wider(names_from = meses_desde, values_from = pct,
              names_prefix = "mes_")
\end{rcode}

\begin{rsalida}
# A tibble: 10 x 2
   cohorte     n
   <chr>   <int>
 1 2023-01    80
 2 2023-02    36
 3 2023-03    29
 4 2023-04    16
 5 2023-05    14
 6 2023-06     3
 7 2023-07     9
 8 2023-08     6
 9 2023-09     3
10 2023-12     2
# A tibble: 3 x 5
  cohorte mes_0 mes_1 mes_2 mes_3
  <chr>   <dbl> <dbl> <dbl> <dbl>
1 2023-01   100    39    38    30
2 2023-02   100    36    19    31
3 2023-03   100    34    38    28
\end{rsalida}

\textbf{Interpretación:} 80 de los 198 clientes activos hicieron su primera
compra en enero, y las cohortes se achican mes a mes porque con compras
aleatorias casi todos los clientes aparecen pronto (solo 2 debutan en
diciembre). En las cohortes del primer trimestre, entre 34\% y 39\% vuelve a
comprar al mes siguiente y alrededor de 30\% al tercer mes, sin una caída
clara. Ojo al interpretar: las cohortes tardías tienen muy pocos clientes
(3 en junio o septiembre) y sus porcentajes serían puro ruido; por eso
mostramos solo las del primer trimestre. En una empresa real las cohortes se
definen con fecha de alta confiable y varios años de historia.
\end{solucion}

\begin{solucion}{p-5}
La regla de Tukey marca como atípico lo que queda fuera de
$[Q_1 - 1.5\,\mathrm{IQR},\ Q_3 + 1.5\,\mathrm{IQR}]$. La segunda regla
compara el precio de cada transacción con la mediana del mismo producto.

\begin{rcode}
library(tidyverse)
transacciones <- read_csv("datasets/transacciones.csv",
                          show_col_types = FALSE)

# Regla (a): Tukey sobre el total de la transacción
q <- quantile(transacciones$total_transaccion, c(0.25, 0.75))
iqr <- unname(q[2] - q[1])
limites <- c(inferior = q[[1]] - 1.5 * iqr, superior = q[[2]] + 1.5 * iqr)
round(limites, 2)
sum(transacciones$total_transaccion < limites["inferior"] |
      transacciones$total_transaccion > limites["superior"])

# Regla (b): precio muy alejado de la mediana del mismo producto
revision <- transacciones %>%
  group_by(producto_id) %>%
  mutate(precio_mediano = median(precio_venta),
         razon = precio_venta / precio_mediano) %>%
  ungroup() %>%
  mutate(atipico_precio = razon > 3 | razon < 1 / 3)

sum(revision$atipico_precio)

# Lista para Auditoría: los casos más extremos primero
revision %>%
  filter(atipico_precio) %>%
  mutate(distancia = abs(log(razon))) %>%
  arrange(desc(distancia)) %>%
  transmute(id = transaccion_id, tienda = tienda_id,
            vendedor = vendedor_id, producto = producto_id,
            precio = precio_venta, mediana = precio_mediano,
            razon = round(razon, 2)) %>%
  head(6)

# ¿Se concentran en algún vendedor?
revision %>% filter(atipico_precio) %>%
  count(vendedor_id, sort = TRUE) %>% head(3)
\end{rcode}

\begin{rsalida}
inferior superior 
-1365.16  3184.89 
[1] 0
[1] 148
# A tibble: 6 x 7
     id tienda vendedor producto precio mediana razon
  <dbl>  <dbl>    <dbl>    <dbl>  <dbl>   <dbl> <dbl>
1   106      9       12       45   20.1    434.  0.05
2   930      6        6       43   24.2    393.  0.06
3   185      3       22       16   22.7    299.  0.08
4   909      1       20       42   22.6    293.  0.08
5   745      4       22       24   20.2    260.  0.08
6   925      2       19        7   26.4    335.  0.08
# A tibble: 3 x 2
  vendedor_id     n
        <dbl> <int>
1          23     9
2          24     9
3           5     8
\end{rsalida}

\textbf{Interpretación:} la regla de Tukey no marca \emph{ninguna}
transacción: el total máximo (\$2,992) está por debajo del límite superior
(\$3,185) y el límite inferior es negativo. Esto pasa porque la cantidad
(1 a 5) y el precio (20 a 600) están acotados. En cambio, comparar el precio
con el de su propio producto marca 148 transacciones (15\%); las más extremas se cobraron a entre 5\% y
8\% del precio mediano del producto, lo cual confirma
el hallazgo de la auditoría: el problema no son los montos, sino precios
inconsistentes para un mismo producto. Los vendedores 23 y 24 concentran 9
casos cada uno; antes de sacar conclusiones sobre personas, compáralo con su
número total de transacciones (43 y 44), porque con 148 casos repartidos en
25 vendedores, 9 por vendedor está dentro de lo esperable.
\end{solucion}

\begin{solucion}{p-6}
Primero medimos si los descuentos se asocian con más unidades en
\archivo{ventas\_retail.csv}; después simulamos diciembre con un descuento
de 10\% y un aumento aleatorio de unidades.

\begin{rcode}
library(tidyverse)
ventas_retail <- read_csv("datasets/ventas_retail.csv",
                          show_col_types = FALSE)
transacciones <- read_csv("datasets/transacciones.csv",
                          show_col_types = FALSE)

# 1. ¿Con descuento se venden más unidades?
ventas_retail %>%
  mutate(con_descuento = descuento_pct > 0) %>%
  group_by(con_descuento) %>%
  summarise(ventas = n(), cantidad_promedio = round(mean(cantidad), 2),
            .groups = "drop")
t.test(cantidad ~ descuento_pct > 0, data = ventas_retail)$p.value

# 2. Simulación de diciembre con 10% de descuento
ingreso_base <- transacciones %>%
  filter(month(fecha) == 12) %>%
  summarise(total = sum(total_transaccion)) %>%
  pull(total)

set.seed(123)
simulacion <- tibble(aumento_unidades = runif(1000, 0, 0.15)) %>%
  mutate(ingreso = ingreso_base * (1 + aumento_unidades) * 0.90,
         diferencia = ingreso - ingreso_base)

simulacion %>%
  summarise(ingreso_base = round(ingreso_base),
            ingreso_promedio = round(mean(ingreso)),
            pct_pierde = round(100 * mean(diferencia < 0), 1),
            peor_caso = round(min(diferencia)),
            mejor_caso = round(max(diferencia)))

# 3. Aumento de unidades necesario para no perder: (1 + x) * 0.9 = 1
round(100 * (1 / 0.90 - 1), 1)
\end{rcode}

\begin{rsalida}
# A tibble: 2 x 3
  con_descuento ventas cantidad_promedio
  <lgl>          <int>             <dbl>
1 FALSE            188              5.61
2 TRUE             177              5.67
[1] 0.8404565
# A tibble: 1 x 5
  ingreso_base ingreso_promedio pct_pierde peor_caso mejor_caso
         <dbl>            <dbl>      <dbl>     <dbl>      <dbl>
1        96394            93225       74.5     -9633       3366
[1] 11.1
\end{rsalida}

\textbf{Interpretación:} en \archivo{ventas\_retail.csv} las ventas con
descuento no tienen más unidades (5.67 contra 5.61, $p = 0.84$): no hay
evidencia de que los descuentos muevan volumen. Con un descuento de 10\% se
necesita vender 11.1\% más unidades solo para igualar el ingreso. Si el
aumento está entre 0\% y 15\%, la campaña pierde ingreso en 74.5\% de las
simulaciones; en promedio el ingreso de diciembre bajaría de \$96,394 a
\$93,225, con una pérdida de hasta \$9,633 y una ganancia máxima de apenas
\$3,366. Y esto es ingreso: con márgenes estimados de 46\%, el efecto en
utilidad sería peor. Recomendación: no lanzar un descuento general; probarlo
primero en unas tiendas con un grupo de control.
\end{solucion}

\begin{solucion}{p-7}
Las \termino{funciones de ventana} (\codigo{min\_rank()}, \codigo{lag()},
\codigo{cumsum()}) calculan un valor por fila usando otras filas del mismo
grupo. Primero completamos la tabla vendedor-mes con \codigo{complete()},
para que un mes sin ventas cuente como cero.

\begin{rcode}
library(tidyverse)
transacciones <- read_csv("datasets/transacciones.csv",
                          show_col_types = FALSE)

ranking <- transacciones %>%
  mutate(mes = month(fecha)) %>%
  count(vendedor_id, mes, wt = total_transaccion, name = "ventas") %>%
  complete(vendedor_id, mes = 1:12, fill = list(ventas = 0)) %>%
  group_by(mes) %>%
  mutate(lugar = min_rank(desc(ventas))) %>%          # ranking del mes
  group_by(vendedor_id) %>%
  arrange(mes, .by_group = TRUE) %>%
  mutate(acumulado = cumsum(ventas),                  # venta acumulada
         cambio_lugar = lag(lugar) - lugar) %>%       # + = subió
  ungroup()

# Vendedores con más meses en el top 5
ranking %>%
  group_by(vendedor_id) %>%
  summarise(meses_top5 = sum(lugar <= 5),
            ventas_anuales = round(max(acumulado)), .groups = "drop") %>%
  arrange(desc(meses_top5), desc(ventas_anuales)) %>%
  head(5)

# Mayor subida de un mes a otro
ranking %>%
  slice_max(cambio_lugar, n = 1) %>%
  select(vendedor_id, mes, lugar, cambio_lugar, ventas) %>%
  mutate(ventas = round(ventas))
\end{rcode}

\begin{rsalida}
# A tibble: 5 x 3
  vendedor_id meses_top5 ventas_anuales
        <dbl>      <int>          <dbl>
1           7          6          50470
2          11          4          50017
3           3          4          48784
4           6          4          47946
5          15          4          42709
# A tibble: 3 x 5
  vendedor_id   mes lugar cambio_lugar ventas
        <dbl> <dbl> <int>        <int>  <dbl>
1           1    10     3           21   5357
2           2     5     3           21   8119
3          17     3     2           21   9971
\end{rsalida}

\textbf{Interpretación:} el vendedor 7, que fue el primero del año
(\$50,470), estuvo en el top 5 en 6 de los 12 meses; ningún otro supera 4
meses. La mayor subida fue de 21 lugares y la lograron tres vendedores (el 1
en octubre, el 2 en mayo y el 17 en marzo): pasar del lugar 23 o 24 al top 3
de un mes a otro indica que el ranking mensual es muy volátil con pocas
transacciones por vendedor (unas 3 o 4 al mes). Para evaluar desempeño
conviene un ranking trimestral o un promedio móvil, no el mes aislado.
\end{solucion}

\begin{solucion}{p-8}
La función arma las tres hojas como una lista con nombre y la escribe con
\codigo{write\_xlsx()} de \paquete{writexl} (cada elemento de la lista es una
hoja). \codigo{walk()} la aplica a las 10 tiendas.

\begin{rcode}
library(tidyverse)
library(writexl)
transacciones <- read_csv("datasets/transacciones.csv",
                          show_col_types = FALSE)
tiendas   <- read_csv("datasets/tiendas.csv", show_col_types = FALSE)
productos <- read_csv("datasets/productos.csv", show_col_types = FALSE)

dir.create("resultados/gerentes", recursive = TRUE, showWarnings = FALSE)

# KPIs de un conjunto de transacciones
kpis_de <- function(datos) {
  datos %>% summarise(ventas = round(sum(total_transaccion)),
                      transacciones = n(),
                      ticket = round(mean(total_transaccion), 2),
                      clientes = n_distinct(cliente_id))
}

# Promedio de la cadena por tienda (para comparar)
cadena <- transacciones %>%
  group_by(tienda_id) %>%
  group_modify(~ kpis_de(.x)) %>%
  ungroup() %>%
  summarise(across(-tienda_id, mean)) %>%
  mutate(ventas = round(ventas), ticket = round(ticket, 2),
         clientes = round(clientes, 1))

paquete_gerente <- function(id) {
  info  <- filter(tiendas, tienda_id == id)
  datos <- filter(transacciones, tienda_id == id)

  resumen <- bind_rows(
    kpis_de(datos) %>% mutate(nivel = info$nombre_tienda),
    cadena %>% mutate(nivel = "Promedio cadena")) %>%
    relocate(nivel)

  mensual <- datos %>%
    group_by(mes) %>%
    summarise(ventas = round(sum(total_transaccion)),
              transacciones = n(), .groups = "drop") %>%
    mutate(crecimiento_pct = round(100 * (ventas / lag(ventas) - 1), 1))

  top <- datos %>%
    group_by(producto_id) %>%
    summarise(unidades = sum(cantidad),
              ventas = round(sum(total_transaccion)), .groups = "drop") %>%
    left_join(select(productos, producto_id, nombre_producto, categoria),
              by = "producto_id") %>%
    slice_max(ventas, n = 10, with_ties = FALSE)

  archivo <- file.path("resultados/gerentes",
                       str_c("tienda_", sprintf("%02d", id), "_",
                             str_replace_all(info$nombre_tienda, " ", "_"),
                             ".xlsx"))
  write_xlsx(list(Resumen = resumen, Mensual = mensual,
                  `Top productos` = top), archivo)
}

walk(tiendas$tienda_id, paquete_gerente)

archivos <- list.files("resultados/gerentes", pattern = "xlsx$")
length(archivos)
head(archivos, 3)

# Comprobación: leer la hoja Resumen de la tienda 8
readxl::read_excel(file.path("resultados/gerentes", archivos[8]),
                   sheet = "Resumen")
\end{rcode}

\begin{rsalida}
[1] 10
[1] "tienda_01_Sucursal_Centro.xlsx" "tienda_02_Sucursal_Norte.xlsx" 
[3] "tienda_03_Sucursal_Sur.xlsx"   
# A tibble: 2 x 5
  nivel           ventas transacciones ticket clientes
  <chr>            <dbl>         <dbl>  <dbl>    <dbl>
1 Sucursal Outlet 114377           106  1079.     82  
2 Promedio cadena  95999           100   959.     79.5
\end{rsalida}

\textbf{Interpretación:} se generaron los 10 archivos, uno por sucursal,
con nombre ordenable (\archivo{tienda\_01\_\ldots}). La hoja
\emph{Resumen} de la Sucursal Outlet muestra que vendió \$114,377 contra un
promedio de \$95,999 por tienda, con un ticket de \$1,079 contra \$959. Ese
formato de ``mi tienda contra la cadena'' es justo lo que un gerente
necesita. Para automatizarlo cada mes, combina esta función con la
programación de tareas del Módulo~\ref{cap:reportes}. Si quieres dar formato
(colores, anchos de columna) usa \paquete{openxlsx}, que no está instalado
en este entorno.
\end{solucion}
